\section{Vérification des solveurs et portée des contrôles}
\label{sec:verification}

\subsection{Tester des propriétés plutôt qu'une supériorité attendue}

Les tests scientifiques interrogent les opérations numériques avant
d'interpréter un classement des méthodes. Une image constante doit
rester identique pour la chaleur et les deux conductances. Zéro
itération doit restituer une copie indépendante de la donnée initiale.
Des entrées finies doivent produire des sorties finies de même
dimension et de type \texttt{float64}, sans modifier le tableau fourni.
Ces contrôles détectent des erreurs de gestion mémoire ou d'indices
que des images visuellement plausibles pourraient masquer.

La conservation est contrôlée sur une grille non carrée et une donnée
aléatoire déterministe. Les tests utilisent une tolérance absolue de
$2\times10^{-13}$ pour la somme et de $2\times10^{-15}$ pour la moyenne.
Ces seuils concernent les petits tableaux de test et les nombres
fixés d'itérations, non une borne universelle pour des tailles
arbitraires. Le principe du maximum est contrôlé à
$2\times10^{-14}$ près, par rapport aux extrema de l'entrée et sans
réduire artificiellement les sorties à $[0,1]$.

Un pixel isolé au bord gauche d'une image $5\times9$ ne doit pas
influencer le bord droit après un seul pas. Ce cas teste directement
l'absence de périodicité ; le caractère non carré aide à repérer une
confusion entre axes. Un autre test utilise $\dx=2$ et $\dy=1$,
pour distinguer explicitement les contributions horizontales et
verticales. Les paramètres inadmissibles sont rejetés : pas spatial
ou temporel non positif, pas temporel hors de la condition suffisante,
nombre d'étapes négatif ou non entier, $K\leq0$, données non finies
et nom de conductance inconnu.

Les valeurs des deux conductances sont contrôlées dans $[0,1]$.
Avec $K=10^9$, une petite image après cinq étapes doit rejoindre
numériquement la chaleur à $2\times10^{-15}$ près en tolérance
absolue. Ce test vérifie le régime $c\simeq1$ dans le code ; il ne
déclare pas qu'une valeur finie de $K$ rend exactement les modèles
identiques. Le bruit et les sorties déterministes sont également
contrôlés à graines fixées. Aucun test n'exige que Perona--Malik
obtienne un PSNR ou un SSIM supérieur à la chaleur.

La suite comporte aussi des tests fonctionnels de configurations,
agrégats, manifestes et modes de vérification. Les 165 tests ont
réussi dans l'exécution finale de phase 4 et dans l'audit indépendant
signalé pour cette phase. Aucun test ni calcul scientifique n'a été
relancé pour la rédaction du présent rapport.

\subsection{Une solution exacte du problème entièrement discret}

Une validation indépendante du solveur de chaleur utilise les modes
propres du Laplacien discret à flux nul. Pour
$0\leq p<N_y$ et $0\leq q<N_x$, définissons
\begin{equation}
 \phi_{pq}[i,j]=
 \cos\left(\frac{\pi p(i+1/2)}{N_y}\right)
 \cos\left(\frac{\pi q(j+1/2)}{N_x}\right).
 \label{eq:discrete-neumann-mode}
\end{equation}
Pour une suite $v_j=\cos(\theta(j+1/2))$, l'identité trigonométrique
$v_{j+1}+v_{j-1}=2\cos(\theta)v_j$ donne
$v_{j+1}-2v_j+v_{j-1}=-4\sin^2(\theta/2)v_j$.
Lorsque $\theta=\pi q/N_x$, le prolongement vérifie
$v_{-1}=v_0$ et $v_{N_x}=v_{N_x-1}$ : la même relation s'applique
au traitement du bord. En combinant les deux directions, on obtient
\begin{equation}
 L\phi_{pq}=\lambda_{pq}\phi_{pq},\qquad
 \lambda_{pq}=
 -\frac{4}{\dy^2}\sin^2\left(\frac{\pi p}{2N_y}\right)
 -\frac{4}{\dx^2}\sin^2\left(\frac{\pi q}{2N_x}\right).
 \label{eq:discrete-neumann-eigenvalue}
\end{equation}
La trajectoire d'Euler correspondante est donc exactement, en
arithmétique réelle,
\begin{equation}
 U^n=(1+\dt\lambda_{pq})^n\phi_{pq}.
 \label{eq:discrete-neumann-amplitude}
\end{equation}
La comparaison utilise une formule analytique indépendante de
l'implémentation des échanges, et non un second appel au même stencil.

Le cas enregistré prend $N_y=17$, $N_x=19$, $p=2$, $q=3$,
$\dy=1{,}1$, $\dx=0{,}8$, $\dt=0{,}2$ et sept pas.
La valeur propre enregistrée est $-0{,}488260677224588$ ; l'erreur
maximale entre le solveur et \eqref{eq:discrete-neumann-amplitude}
vaut $3{,}608224830031759\times10^{-16}$, très inférieure au seuil
de test $5\times10^{-13}$. Ces valeurs proviennent du bloc
\texttt{heat\_discrete\_neumann\_mode} de
\texttt{results/extended/verification.json}. Elles valident cette
réalisation du Laplacien et d'Euler pour un mode non constant,
notamment avec une grille non carrée et des pas différents.
Elles ne mesurent pas encore une erreur par rapport à une EDP continue.

\subsection{Deux contrôles d'amplitude séparant espace et temps}

Le contrôle spatial considère le mode cosinus $(p,q)=(1,1)$ sur le
carré unité, avec $\dx=\dy=1/N$ et $T=0{,}01$.
L'amplitude continue exacte est $\exp(T\lambda)$, avec
$\lambda=-2\pi^2$. On lui compare l'amplitude semi-discrète exacte
$\exp(T\lambda_h)$ : il n'y a donc aucune erreur d'Euler dans cette
comparaison. Les erreurs enregistrées sont
$5{,}200695560823032\times10^{-4}$ pour $N=16$ et
$1{,}3011185809019832\times10^{-4}$ pour $N=32$.

Le contrôle temporel fixe au contraire la grille $N_y=16$, $N_x=20$
sur le carré unité, les pas $\dy=1/16$, $\dx=1/20$, le mode
$(p,q)=(1,2)$ et $T=0{,}001$. Il compare
$(1+(T/m)\lambda_h)^m$ à $\exp(T\lambda_h)$, pour $m=8$ puis
$m=16$ pas. Les erreurs sont respectivement
$1{,}4342037335590696\times10^{-4}$ et
$7{,}156608836045297\times10^{-5}$.
Ces quatre erreurs, lues dans le bloc \texttt{refinement\_checks}
du même compte rendu, diminuent dans les deux comparaisons séparées.
Il s'agit de contrôles d'amplitude fondés sur les expressions
spectrales, non d'un nouveau benchmark d'images.

Deux niveaux ne constituent pas une preuve expérimentale d'un ordre
général de convergence. De plus, les modes cosinus sont très réguliers
et particulièrement compatibles avec le bord de Neumann. Ils ne
représentent pas les discontinuités géométriques et le bruit des images
du protocole. Le rapport se limite donc au constat de diminution
mesurée dans ces contrôles précis.

\subsection{Consistance : intérieur, bord et régularité}

Pour la chaleur, sur une fonction suffisamment régulière à l'intérieur,
les différences centrées donnent une erreur spatiale
$O(\dx^2+\dy^2)$ par développement de Taylor. Euler explicite possède
une erreur de consistance temporelle $O(\dt)$ lorsque la solution
dispose des dérivées temporelles nécessaires. Ces énoncés locaux ne
prouvent pas, à eux seuls, la convergence globale pour toute donnée.

Le bord mérite une réserve explicite. En une dimension, au premier
centre $x=h/2$, le flux nul impose
$(u(3h/2)-u(h/2))/h^2$ comme contribution du Laplacien.
Pour $u'(0)=0$ et une fonction assez régulière, son développement est
\begin{equation}
 \frac{u(3h/2)-u(h/2)}{h^2}
 =u''(h/2)+\frac{h}{24}u'''(0)+O(h^2).
 \label{eq:boundary-consistency}
\end{equation}
La précision ponctuelle générique est donc seulement du premier
ordre au bord sous la seule condition de Neumann. Si la compatibilité
supplémentaire $u'''(0)=0$ est satisfaite, ce terme disparaît ; c'est
le cas des modes cosinus utilisés. Un énoncé d'ordre deux sur tout
le rectangle, sans mention de ces hypothèses ou d'une analyse globale
adaptée, serait excessif.

Enfin, pour la variante Perona--Malik, les développements intérieurs
sur des fonctions régulières doivent être rapportés à l'opérateur
directionnel \eqref{eq:directional-limit}. Ils ne peuvent pas être
utilisés pour revendiquer une consistance isotrope avec
\eqref{eq:pm-continuous}. Aux frontières et pour des images bruitées
ou discontinues, les mêmes exigences de régularité doivent être
réexaminées. La vérification du code et la preuve de ses propriétés
discrètes restent ainsi distinguées d'une théorie complète du
problème continu.
